\documentclass[aspectratio=169,11pt]{beamer}
\usepackage{rankit}

\usetheme{Boadilla}
\setbeamercolor{structure}{fg=RankItTeal}
\setbeamercolor{frametitle}{fg=white,bg=RankItNavy}
\setbeamercolor{title}{fg=RankItNavy}
\setbeamercolor{block title}{fg=white,bg=RankItNavy}
\setbeamercolor{block body}{fg=RankItInk,bg=RankItPale}
\setbeamertemplate{navigation symbols}{}
\setbeamertemplate{footline}[frame number]
\hypersetup{pdftitle={RankIt - repozytorium i szablon},pdfauthor={Zespół RankIt}}

\title{RankIt}
\subtitle{Repozytorium i szablon dokumentów}
\author{Igor Jankowski \and Piotr Ratkowski \and Zachariasz Słupski}
\institute{Realizacja Projektu Informatycznego 2026/2027}
\date{24 października 2026}

\begin{document}

\begin{frame}[plain]
  \titlepage
\end{frame}

\begin{frame}{Projekt RankIt}
  \begin{block}{Cel aplikacji}
    Prosta aplikacja webowa do tworzenia rankingów i oceniania różnych rzeczy.
  \end{block}
  \begin{itemize}
    \item Użytkownik tworzy ranking i dodaje do niego pozycje.
    \item Oceniający przyznają pozycjom oceny.
    \item Wyniki są uporządkowane i łatwe do porównania.
  \end{itemize}
  \vfill
  \small Zakres pierwszej wersji jest propozycją do zatwierdzenia przez zespół.
\end{frame}

\begin{frame}{Repozytorium projektu}
  \textbf{Adres:} \url{https://github.com/elprofesoriqo/rpi}
  \vspace{0.6cm}
  \begin{itemize}
    \item Publiczne repozytorium przechowuje kod, dokumenty i historię zmian.
    \item \texttt{main} jest gałęzią główną.
    \item Zadania zapisujemy w GitHub Issues.
    \item Większe zmiany sprawdzamy w Pull Requestach.
  \end{itemize}
\end{frame}

\begin{frame}{Skład zespołu}
  \begin{itemize}
    \item Igor Jankowski - koordynacja, repozytorium i dokumentacja.
    \item Piotr Ratkowski - projekt widoków i interfejs użytkownika.
    \item Zachariasz Słupski - logika rankingu, dane i scenariusze sprawdzeń.
  \end{itemize}
  \small Podział zadań jest propozycją organizacyjną, nie oceną wcześniejszych kompetencji.
\end{frame}

\begin{frame}{Komunikacja}
  \begin{itemize}
    \item Messenger służy do bieżącej komunikacji zespołu.
    \item GitHub Issues przechowują zadania i decyzje dotyczące kodu.
    \item Proponowane jest krótkie spotkanie zespołu raz w tygodniu.
    \item Pytania do opiekuna kierujemy kanałem wskazanym przez prowadzącego.
  \end{itemize}
\end{frame}

\begin{frame}{Szablon dokumentu}
  Każdy dokument zawiera:
  \begin{itemize}
    \item projekt, tytuł i identyfikator dokumentu;
    \item wersję, datę, autorów, status i opiekuna;
    \item historię zmian, treść, wnioski i źródła;
    \item stałą metrykę w nagłówku i stopce.
  \end{itemize}
  \vfill
  \begin{block}{Identyfikacja wizualna}
    Granat \texttt{\#243B53}, turkus \texttt{\#2CB1BC} i znak \RankItWordmark.
  \end{block}
\end{frame}

\begin{frame}{Narzędzia i pliki}
  \begin{itemize}
    \item Git i GitHub - kod, przegląd zmian i historia projektu.
    \item Messenger - bieżące wiadomości zespołu.
    \item LaTeX - edytowalne źródła raportu, szablonu i prezentacji.
    \item PDF - gotowe materiały do przesłania na kurs.
  \end{itemize}
  \vfill
  \small Kod aplikacji i stos technologiczny nie zostały jeszcze określone.
\end{frame}

\begin{frame}{Wykorzystanie AI}
  \begin{itemize}
    \item AI pomogło uporządkować wymagania i przygotować pierwszą wersję materiałów.
    \item Zespół przekazał skład i rzeczywisty zrzut grupy Messenger.
    \item Propozycje zakresu, podział zadań i spotkań wymagają zatwierdzenia przez zespół.
    \item Nazwisko opiekuna i termin oddania trzeba uzupełnić po uzyskaniu danych.
  \end{itemize}
\end{frame}

\end{document}
